% No preamble. Notebook 15 candidate module, ot- labels.
\section{Transfer to the ordinary homogeneous averaging constraint}
\label{ot-section}

The preceding minimizer was constructed for the corrected
fixed-reference constraint $C^{(j)}(U)=D=e^{iB}$. The ordinary
constraint is $M^{(j)}(U)=D$, with $M^{(j)}$ the homogeneous
iteration of the ordinary group average. These two equalities
are distinct. This section proves a transfer between their small
field minimizers using actual gauge transformations. All fields
are on finite compatible periodic regulators; final spacing is
one. The radii and constants below are independent of refinement
and volume. The bare Wilson action is gauge invariant.

\subsection{An exact multilevel gauge identity}

Let $C_t=C^{(t)}(U)$ and $b_0(U)=1$. Notebook 14 proves
\begin{equation}\label{ot-cumulative}
 b_{t+1}(U)=b_t(U)|_{t+1}\,a(C_t),\qquad
 M^{(t)}(U)=C_t^{\,b_t(U)}.
\end{equation}
Here restriction is to the next coarse sites, and
$a(C_t)$ is the endpoint mean at that step. In particular,
on a fixed sufficiently small bounded-potential ball,
\begin{equation}\label{ot-b-bound}
 \norm{b_j(U)-1}_\infty+\norm{b_j(U)^{-1}-1}_\infty
       \le C\norm{A}_\infty,\qquad U=e^{i\eta A}.
\end{equation}
These group-valued maps are holomorphic with uniform bounds
on a common complex ball. This follows either directly from
the product in \eqref{ot-cumulative} and the summability of
the physical spacings, or from the established corrected
trajectory bounds. Only a coarse supremum bound is needed here.

For a fine gauge $g$, let $k_0=g$ and let $k_t$ be its induced
gauge under the corrected iteration, so that
$C^{(t)}(U^g)=C_t^{k_t}$. The one-step formula gives
\[
 k_{t+1}=a(C_t^{k_t})^{-1}\,k_t|_{t+1}\,a(C_t).
\]
Multiplication with \eqref{ot-cumulative}, applied also to
$U^g$, proves by induction
\begin{equation}\label{ot-exact-induced}
 b_t(U^g)k_t=g|_t\,b_t(U),\qquad
 k_j(U,g)=b_j(U^g)^{-1}\,g|_j\,b_j(U).
\end{equation}
Indeed the adjacent factors $a(C_t^{k_t})$ and their inverses
cancel before the inductive identity is used. No multiplicative
cocycle for the logarithmic endpoint average is assumed.

\subsection{Uniform analytic action of smooth fine gauges}

Write $\mathcal X_\alpha=C^{1,\alpha}$ for the complete
potential norm and $\mathcal Z_\alpha=C^{2,\alpha}$ for scalar
gauge parameters, as in the preceding section. There is a
uniform analytic map
\begin{equation}\label{ot-gauge-action}
 \mathscr A_\eta(A,\lambda)(e)
  =\frac1{i\eta}\log\bigl[
      e^{i\lambda(e_-)}e^{i\eta A(e)}e^{-i\lambda(e_+)}\bigr]
\end{equation}
from a fixed small ball in
$\mathcal X_\alpha\times\mathcal Z_\alpha$ to
$\mathcal X_\alpha$, and
\begin{equation}\label{ot-gauge-action-bound}
 \norm{\mathscr A_\eta(A,\lambda)}_{\mathcal X_\alpha}
       \le C\bigl(\norm{A}_{\mathcal X_\alpha}
                       +\norm{\lambda}_{\mathcal Z_\alpha}\bigr).
\end{equation}
To justify the apparent division by $\eta$, put
$\ell=\lambda(e_-)$ and $v=\nabla_\mu\lambda(e_-)$.
The group element in \eqref{ot-gauge-action} equals
\[
 e^{i\eta\operatorname{Ad}_{e^{i\ell}}A(e)}
        e^{i\ell}e^{-i(\ell+\eta v)}.
\]
For every small $(A,\ell,v)$ this analytic product is one
at $\eta=0$. Its logarithm is therefore divisible by $\eta$
as an analytic function, with a uniform extension for
$0\le\eta\le1$. All its finite-dimensional derivatives are
uniform on a smaller fixed complex jet ball. Discrete product,
shift and analytic-composition estimates applied to
$(A,\lambda,\nabla\lambda)$ prove
\eqref{ot-gauge-action-bound} and holomorphy in the indicated
Banach norms. A first potential difference uses at most two
scalar differences. The same argument applies to its
$\alpha$-H\"older seminorm, so no inverse-mesh loss occurs.
For real Hermitian traceless arguments the output remains
Hermitian traceless on the common logarithm branch.

\subsection{Prescribing the induced coarse gauge}

Use the smooth block-mean right inverse $H_b$ constructed in
\eqref{ng-biharmonic-extension}--\eqref{ng-biharmonic-bounds}.
It is a scalar operator tensored with the gauge Lie algebra.
For bounded coarse $\theta$ define
\begin{equation}\label{ot-K-definition}
 \mathcal K(A,\theta)=\frac1i\log
      k_j\bigl(e^{i\eta A},e^{iH_b\theta}\bigr).
\end{equation}
The exact formula \eqref{ot-exact-induced}, the bound
\eqref{ot-b-bound}, the gauge-action estimate and
$H_b:\ell^\infty\to\mathcal Z_\alpha$ show that
$\mathcal K$ is uniformly bounded and holomorphic on a
fixed ball in $\mathcal X_\alpha\times\ell^\infty$.
Evaluation of $H_b\theta$ at coarse sites is bounded by
its supremum norm. Thus this step does not presume that
$H_b$ interpolates point values.

At $\theta=0$ the induced gauge is one, so
$\mathcal K(A,0)=0$. At the joint origin its differential is
\begin{equation}\label{ot-K-linear}
 D_\theta\mathcal K(0,0)=\mathcal G_0H_b=Q_0H_b=I.
\end{equation}
Uniform Cauchy estimates on a smaller ball consequently give
\begin{equation}\label{ot-K-near-identity}
 \norm{D_\theta\mathcal K(A,\theta)-I}_{\infty\to\infty}
       \le C\bigl(\norm{A}_{\mathcal X_\alpha}
                                    +\norm{\theta}_\infty\bigr).
\end{equation}
For small $A$ and any sufficiently small bounded coarse
$\omega$, the contraction
\[
 \theta=\omega-\bigl[\mathcal K(A,\theta)-\theta\bigr]
\]
therefore has a unique solution in a fixed small ball.
The solution satisfies $\norm{\theta}_\infty\le2\norm{\omega}_\infty$
and depends holomorphically on $(A,\omega)$, with uniform
Cauchy bounds on smaller balls. It is real for real inputs.
Thus any small prescribed coarse gauge $e^{i\omega}$ is
realized as the induced corrected gauge of a fine gauge
$e^{iH_b\theta}$ of uniformly controlled smoothness.

Suppose now that $U=e^{i\eta A}$ satisfies
$M^{(j)}(U)=D$. By \eqref{ot-cumulative},
$C^{(j)}(U)^{b_j(U)}=D$. Prescribe
\[
 \omega=\frac1i\log b_j(U),\qquad k_j(U,g)=b_j(U).
\]
The preceding construction produces a fine gauge $g$ such that
\begin{equation}\label{ot-ordinary-to-corrected}
 C^{(j)}(U^g)=D,\qquad
 \mathcal S_\eta(U^g)=\mathcal S_\eta(U),\qquad
 \norm{\mathscr A_\eta(A,\log(g)/i)}_{\mathcal X_\alpha}
                     \le C\norm{A}_{\mathcal X_\alpha}.
\end{equation}
Here $\log(g)/i=H_b\theta$ on the chosen branch. This map
takes sufficiently small ordinary-feasible fields into
uniformly small corrected-feasible fields.

\subsection{An ordinary-feasible representative of the minimizer}

Let $U_*(B)=e^{i\eta A_*(B)}$ be the corrected minimizer
constructed above, so $C^{(j)}(U_*)=D=e^{iB}$ and
$\norm{A_*(B)}_{\mathcal X_\alpha}\le C\norm{B}_\infty$.
We need a smooth fine gauge with specified \emph{point values}
at the final coarse sites. For this purpose construct a
different operator $E_{\rm pt}$ as follows. Choose a fixed
smooth scalar bump $\varphi$ supported in
$(-1/2,1/2)^d$ and satisfying $\varphi(0)=1$. On the
periodic coarse lattice put
\begin{equation}\label{ot-point-extension}
 (E_{\rm pt}\omega)(x)=\sum_{n\in\mathbb Z^d}
                  \varphi(x-n)\omega(n),
\end{equation}
with the coefficients extended periodically; evaluate the
result at fine sites. This is well defined on every
compatible torus, including the smallest coarse periods.
Its derivatives through order three are bounded by
$C\norm{\omega}_\infty$, since the translates have uniformly
bounded overlap. Sampling and the integral formula for
finite differences give, uniformly in $\eta$,
\begin{equation}\label{ot-point-bounds}
 E_{\rm pt}\omega|_{\rm coarse}=\omega,\qquad
 \norm{E_{\rm pt}\omega}_{\mathcal Z_\alpha}
                          \le C\norm{\omega}_\infty.
\end{equation}
Set
\begin{equation}\label{ot-ordinary-representative}
 \omega_*(B)=\frac1i\log b_j(U_*(B))^{-1},\quad
 h_*(B)=e^{iE_{\rm pt}\omega_*(B)},\quad
 U_{\rm ord}(B)=U_*(B)^{h_*(B)}.
\end{equation}
Ordinary gauge covariance and $h_*|_{\rm coarse}=b_j(U_*)^{-1}$
give
\begin{equation}\label{ot-ordinary-exact}
 M^{(j)}(U_{\rm ord}(B))=D,\qquad
 \mathcal S_\eta(U_{\rm ord}(B))=\mathcal S_\eta(U_*(B)).
\end{equation}
Indeed with the convention $V^g=g_-Vg_+^{-1}$,
$(V^b)^h=V^{hb}$, so the two coarse gauges cancel in the
required order. The ordinary potential
$A_{\rm ord}(B)=\log U_{\rm ord}(B)/(i\eta)$ is holomorphic
on a uniform coarse supremum ball and obeys
\begin{equation}\label{ot-ordinary-bound}
 \norm{A_{\rm ord}(B)}_{\mathcal X_\alpha}
                              \le C\norm{B}_\infty.
\end{equation}

\begin{theorem}[Ordinary homogeneous constrained minimum]
\label{ot-minimum}
Fix $0<\alpha<1$, a fixed odd blocking ratio $L\ge3$,
and the compact matrix gauge group and trace normalization
used in the Wilson calculation. There are $b_0,r_0>0$,
independent of the homogeneous refinement $M=L^j$ and
finite compatible periodic volume, with the following
property. For every real coarse potential
$\norm{B}_\infty<b_0$, the field $U_{\rm ord}(B)$ in
\eqref{ot-ordinary-representative} satisfies the ordinary
constraint $M^{(j)}(U)=e^{iB}$ and minimizes the bare Wilson
action among all fields $U=e^{i\eta A}$ satisfying that
constraint and $\norm{A}_{\mathcal X_\alpha}<r_0$.
Equality in this small field class implies gauge equivalence
to $U_{\rm ord}(B)$. Its ordinary potential has the uniform
bound \eqref{ot-ordinary-bound} and a holomorphic extension
to a fixed complex coarse-data ball.
\end{theorem}

\begin{proof}
The reduced-Hessian proof gives strict convexity of the
corrected fixed-slice action on a fixed ball of the linear
coordinate $\xi\in X$, for all sufficiently small $B$.
In particular the corrected minimizer is the unique minimum
on that ball. The constraint chart covers every sufficiently
small corrected-feasible fixed-slice field, and
$\xi=(I-HQ)A$ there; this coordinate is uniformly small in
$\mathcal X_\alpha$ since $Q$ is bounded into coarse
supremum norm and $H$ maps that norm into
$\mathcal X_\alpha$.

Given an ordinary-feasible field with
$\norm{A}_{\mathcal X_\alpha}<r_0$, first apply
\eqref{ot-ordinary-to-corrected}, and then the
constraint-preserving flow into the flat slice. Both
operations preserve the action and enlarge the potential
norm by at most a fixed factor. Choose $r_0$ so small that
the resulting fixed-slice field belongs to the chart and
its $\xi$ belongs to the strict-convexity ball. The action
comparison with $U_*(B)$ now follows from that convexity.
Choose $b_0$ smaller, if necessary, so that
$\norm{A_{\rm ord}(B)}_{\mathcal X_\alpha}<r_0$ and the
minimizer coordinate belongs to the same convex ball.
Equation \eqref{ot-ordinary-exact} transfers the comparison
to $U_{\rm ord}(B)$. All radii can be chosen with fixed
strict margins and are uniform in refinement and volume.

If equality holds, strict convexity identifies the final
fixed-slice field with $U_*(B)$. Both earlier operations
are genuine gauge transformations, so the original field
is gauge equivalent to $U_{\rm ord}(B)$. Since both have
the same ordinary coarse constraint, the restriction of
a connecting gauge to the coarse sites stabilizes that
coarse field. No uniqueness of the gauge parameter is
claimed. The analyticity and norm statements follow from
\eqref{ot-b-bound}, \eqref{ot-gauge-action-bound} and the
analytic corrected minimizer.
\end{proof}

This is a minimum in a fixed small \emph{global potential}
class on each finite homogeneous regulator. It is not a
global minimum over all lattice gauge fields, nor a claim
about arbitrary local-frame or mixed-history constraints.
The source's notation $M_{\mathfrak B}(U)=D$ includes more
general geometries; only the explicitly defined ordinary
homogeneous iteration is proved here. No continuum quantum
measure or physical mass gap follows from this theorem.
